\documentclass[10pt,a4paper]{amsart}
\usepackage[margin=19mm]{geometry}
\usepackage{amsmath,amssymb,mathtools}
\usepackage[scr=rsfs,cal=euler]{mathalfa}
\usepackage{xfrac}
\usepackage[unicode,hidelinks,bookmarks=false]{hyperref}
\newcommand{\ie}{i.e.\ }
\newcommand{\cf}{cf.\ }
\newcommand{\resp}{respectively\ }
\newcommand{\Subsection}{\subsection}
\providecommand{\nobreakdash}{\nobreak\hbox{-}}
\setlength{\emergencystretch}{2em}
\multlinegap0pt
\usepackage{bm}
\usepackage{bbm}
\usepackage{dsfont}
\usepackage{xspace}
\usepackage{cancel}
\usepackage{mathtools}
\usepackage{amsthm}
\makeatletter
\@ifundefined{lemma}{\newtheorem{lemma}{Lemma}}{}
\@ifundefined{proposition}{\newtheorem{proposition}{Proposition}}{}
\makeatother
\title[Non-Existence of Linear-Quartic Factorization for the Second]{Non-Existence of Linear-Quartic Factorization for the Second Cuboid Quintic}
\author{Valery Asiryan}
\date{Source: arXiv:2601.04241v2 (CC BY 4.0)}
\begin{document}
\maketitle

\noindent\emph{Source and licence.} Non-Existence of Linear-Quartic Factorization for the Second Cuboid Quintic. Version arXiv:2601.04241v2 (CC BY 4.0), distributed under the Creative Commons Attribution 4.0 International licence, as recorded in the arXiv licence metadata for this exact version and shown on the abstract page https://arxiv.org/abs/2601.04241v2. Full rights notice: licenses/arxiv-2601.04241-CC-BY-4.0.txt.\par\smallskip

\noindent\textit{Editorial scope.} Counted unit: the Lemma whose source label is \texttt{lem:noadmissible} in the article (source0.tex lines 393--396, source label \texttt{lem:noadmissible}), delivered together with the article's own complete original proof of it (source0.tex lines 398--422, one \texttt{proof} environment of 25 lines). No mathematics is written, repaired or re-derived: statement and proof are the article's own text line for line. Required context retained uncounted: eq:Uparam (lines 294--296); eq:Ccurve (lines 340--342); prop:CQ (lines 379--385). Every definition, display and label of those lines is retained unchanged. Omitted from this package and not redistributed: 1--293, 297--339, 343--378, 386--392, 397--397, 423--597. Nothing inside a retained excerpt is omitted or altered: every retained body is delivered exactly as the article's lines are written, so no omission receipt of any kind accompanies this package. Citations: the retained excerpts carry no citation command at all, so no bibliography entry of the article is transcribed and no reference is redistributed. Reference adaptation: none was needed and none was made; every reference inside the delivered unit resolves inside it, so no unresolved reference or citation is produced. Printed slips kept literal: no printed word, symbol, hypothesis or number of the counted statement or of its proof was altered. Layout and class: the article's own class is \texttt{article} with packages including inputenc, fontenc, amsmath, amssymb, amsthm, geometry, hyperref, indentfirst, booktabs, xcolor, fvextra; the wrapper is the lane's pinned \texttt{amsart} editorial wrapper at 10pt on A4, which restates the theorem-like environments the retained text needs and adds the article's own macro definitions that the retained text needs (none), with extra wrapper packages (bm, bbm, dsfont, xspace, cancel, mathtools). The delivered numbering is the unit's own. Assets: no retained span contains a figure environment, a graphics-inclusion command, a PDF-page inclusion, a table or a TikZ picture; no image file is copied into this package, the package declares no asset dependency and no image file is redistributed with it. Licence: version arXiv:2601.04241v2 of this article is distributed under the Creative Commons Attribution 4.0 International licence, as recorded in the arXiv licence metadata for this exact version and shown on the abstract page https://arxiv.org/abs/2601.04241v2. A full rights notice is delivered at licenses/arxiv-2601.04241-CC-BY-4.0.txt. Characters: every retained excerpt is pure ASCII, so the delivered engine, XeLaTeX with its default Latin Modern text font, reports no missing character.
\begin{equation}\label{eq:Uparam}
U(t)=\frac{2\,(t^{5}+6t^{4}+30t^{3}+16t^{2}+9t+2)}{t\,(t-1)^{2}(t^{2}+6t+1)},
\end{equation}

\begin{equation}\label{eq:Ccurve}
C:\quad w^{2}=(t+1)(t^{4}+20t^{3}+6t^{2}+4t+1).
\end{equation}

\begin{proposition}[Rational points on $C$]\label{prop:CQ}
Let $C$ be the hyperelliptic curve \eqref{eq:Ccurve}.
Then
\[
C(\mathbb{Q})=\bigl\{\infty,\ (-1,0),\ (0,\pm 1),\ (1,\pm 8)\bigr\}.
\]
\end{proposition}

\begin{lemma}[No admissible parameters from $C(\mathbb{Q})$]\label{lem:noadmissible}
Let $(t,w)\in C(\mathbb{Q})$ and let $U=U(t)$ be as in \eqref{eq:Uparam} whenever defined.
Then the only positive rational value of $s$ satisfying $U=s+1/s$ is $s=1$.
\end{lemma}

\begin{proof}
By Proposition~\ref{prop:CQ}, the rational points on $C$ correspond to $t \in \{-1, 0, 1, \infty\}$. We determine the admissible values of $s$ by solving $s+1/s = U(t)$ using the expression \eqref{eq:Uparam}.

\smallskip
\noindent\emph{Case $t=\infty$.}
Comparing the degrees of the numerator and denominator in \eqref{eq:Uparam}, we find $\lim_{t\to\infty} U(t) = 2$. The equation $s+1/s=2$ is equivalent to $(s-1)^2=0$, yielding the unique solution $s=1$. This corresponds to the excluded case $p=q$.

\smallskip
\noindent\emph{Case $t=-1$.}
Substituting $t=-1$ into \eqref{eq:Uparam} yields
\[
U(-1) = \frac{2\,(-1+6-30+16-9+2)}{(-1)\,(-2)^{2}\,(1-6+1)} = \frac{-32}{16} = -2.
\]
The equation $s+1/s=-2$ implies $(s+1)^2=0$, so $s=-1$. However, the cuboid parameter $s=(p/q)^2$ must be positive, so this solution is inadmissible.

\smallskip
\noindent\emph{Cases $t=0$ and $t=1$.}
The denominator of $U(t)$ in \eqref{eq:Uparam} contains the factors $t$ and $(t-1)^2$, so it vanishes at $t=0$ and $t=1$. Evaluating the numerator at these points yields $4$ (for $t=0$) and $128$ (for $t=1$), which are non-zero.
Consequently, $t=0$ and $t=1$ are poles of the rational function $U(t)$, implying $|U(t)| \to \infty$.
The relation $s+1/s=U$ implies that $s \to 0$ or $s \to \infty$. Since we require $s \in \mathbb{Q}_{>0}$ (a finite non-zero rational), these points yield no admissible parameters.

\smallskip
\noindent\emph{Conclusion.}
The only positive rational $s$ arising from $C(\mathbb{Q})$ is $s=1$.
\end{proof}

\end{document}
